\section{Related Work} \label{chapter:RELATED_WORK}

Related work is two-fold: Approaches related to recovering previous system states and approaches related to our technical solution and, thus, to scoping and grouping changes into first-class objects.

\subsection{Recovering Previous System States}

\emph{\coexist}\cite{Steinert2012COE} provides recovery support through continuous versioning in Squeak/Smalltalk.
For each change made to source code, \coexist creates a new version of the system sources, resulting in a fine-grained history of changes without requiring precautionary actions by the developer.
As in our approach, preserved versions are part of the program runtime and can be re-established easily.
Similar to this approach are \emph{Changeboxes}\cite{Denker2007EEC}, an approach to capturing and scoping changes to a system using first-class entities called Changeboxes.
However, Changeboxes are primarily useful for reviewing the evolution of a system. 
In contrast, to undo changes using Changeboxes is rather tedious~\cite{Steinert2012COE}.
However, while both \coexist and ChangeBoxes preserve only changes to
the source code of classes, our system preserves also the state and
behavior of objects.

\emph{Back-in-time Debuggers}~\cite{Lewis2003BIT}, also known as \emph{Omniscient Debuggers}, allow developers to inspect previous program states and step backwards in the control flow to undo the side effects of statements.
Approaches for this are either based on logging or replay: either the debugger records information to be able to recreate particular previous situations, requiring mainly space for the different states, or the debugger re-executes the program up to a particular previous situation, requiring mainly time to re-run the program.
Our approach is similar to logging-based back-in-time debugging, re-establishing a previous state through preserving information.
However, back-in-time debuggers need to be able to undo the effects of each statement separately, while our system's versioning granularity is arbitrarily and can, for example, correspond to programmer interactions with the system.
Additionally, we intend object versioning to be active during all development tasks, not only when debugging.

\emph{\Acp{stm}}~\cite{Shavit1995STM} captures changes to values in transactions, analogous to database transactions.
Each transaction has its own view of the memory, which is unaffected by other concurrently running transactions.
Multiple versions of the system state can coexist and which version is read and written to depends on the transaction.
\ac{stm} and our approach are similar in that multiple versions of the system state can coexist and that a previous state can be re-established if necessary.
However, \ac{stm} provides concurrency control and an alternative to lock-based synchronization, while our approach provides recovery support to developers when changes turn out be inappropriate.
Due to these different goals, our versions are also first-class objects, which can be stored in variables and be re-established at any time, while \ac{stm} transactions are created implicitly through particular control structures and committed immediately upon success.

\subsection{Dynamically Scoping First-class Groups of Changes}

\emph{Worlds}\cite{Warth2011Wor} provide a language construct for controlling the scope of side effects: changes to the state of objects are by default only effective in the \emph{world} in which the changes occurred.
These worlds are first-class values can be spawned from an existing world, which establishes a child-parent relationship between the two worlds.
Worlds provide a language construct for experimenting with different states of the system, while object versioning allows to preserve versions of the system to recover previous states: Our approach does not include extensions to the host programming languages and no conditions for combining versions with their predecessor versions.

\emph{Object Graph Versioning}\cite{Pluquet2009ECP} allows programmers to preserve access to previous states of objects.
Fields of objects can be marked as \emph{selected} fields.
When a snapshot is created, the values of these selected fields are preserved.
Therefore, not every state can be re-established, but states that are part of global snapshots. 
The approach, thus, provides fine-grained control to programmers regarding which fields of which objects should be preserved when.
Since Object Graph Versioning aims to support implementing application-specific undo/redo or tools, individual fields are versioned only when programmers explicitly mark them as selected.

\emph{Practical Object-oriented Back-in-Time Debugging}\cite{Lienhard2008POB} is a logging-based approach to back-in-time debugging that uses alternative references to preserve the history of objects.
These alternative references, called \emph{Aliases}, are actually objects and part of the application memory.
They contain information about the history and origin of the values stored in fields.
Aliases can not only revert to previous states of the system, but also retrace the flow of all values, while object-versioning only preserves the system states.
However, as with back-in-time debugging, the alias references are intended to be used in explicit debugging sessions, while our version-aware reference are intended to be used at all times.

%%% Local Variables:
%%% mode: latex
%%% TeX-master: "../main"
%%% End:
